% Generated draft: requires mathematical and source review.
\documentclass[11pt]{article}
\usepackage[margin=25mm]{geometry}
\usepackage{fontspec}
\setmainfont{Latin Modern Roman}
\usepackage{amsmath,amssymb,mathtools,unicode-math}
\setmathfont{Latin Modern Math}
\usepackage{enumitem,needspace,xurl,hyperref}
\setlength{\parindent}{0pt}
\setlength{\parskip}{5pt}
\setlength{\emergencystretch}{3em}
\newcommand{\N}{\mathbb{N}}
\newcommand{\Z}{\mathbb{Z}}
\newcommand{\Q}{\mathbb{Q}}
\newcommand{\R}{\mathbb{R}}
\newcommand{\C}{\mathbb{C}}
\newcommand{\F}{\mathbb{F}}
\newcommand{\A}{\mathbb{A}}
\newcommand{\PP}{\mathbb{P}}
\newcommand{\E}{\mathbb{E}}
\newcommand{\eps}{\varepsilon}
\newcommand{\dd}{\,\mathrm{d}}
\begin{document}
\section*{UnsolvedMath ID: 20002090}
\textbf{Source problem code:} catalogue entry
\subsection{Definitions and mathematical statement}
\textbf{Source scope.} The complete Question 1.32 says only that Akahori discussed problems about the geometry of Kohn--Rossi cohomology. It supplies no manifold class, bidegree, conjecture, estimate or requested construction. It is therefore context only; the following regular-circle-bundle calculation is an auxiliary example from the aid, not a reconstructed missing question.

Let $B$ be a compact complex manifold of dimension $n\ge1$ and let $L\to B$ be a positive holomorphic line bundle with a positively curved Hermitian metric. Its dual unit circle bundle $X=S(L^*)$ has its standard regular strongly pseudoconvex CR structure and circle action $R_\theta(v)=e^{i\theta}v$. The scalar Kohn--Rossi groups are the cohomology of the tangential Cauchy--Riemann complex,
\[
 H_b^{0,q}(X)=\ker(\bar\partial_b:\Omega_b^{0,q}\to\Omega_b^{0,q+1})
          /\operatorname{im}(\bar\partial_b:\Omega_b^{0,q-1}\to\Omega_b^{0,q}).
\]
Use weight $m\in\Z$ to mean $R_\theta^*u=e^{im\theta}u$. With this convention the aid asserts the natural Fourier--Dolbeault identification
\[
 H_b^{0,q}(X)_m\cong H^q(B,L^m),\qquad 0\le q\le n.
\]
In particular positive weight corresponds to $L^m$, not $L^{-m}$. Reversing the circle action reverses that sign.

For each interior degree $1\le q\le n-1$, the aid asserts that only finitely many weights contribute, giving
\[
 H_b^{0,q}(X)\cong\bigoplus_{m\in\Z}H^q(B,L^m),
\]
with a finite nonzero sum. The positive tail is controlled by Serre vanishing and the negative tail by Serre duality and vanishing for $K_B\otimes L^{-m}$; no effective uniform numerical cutoff is supplied. When $n=1$ the interior range is empty. The finite-window assertion does not extend to endpoint degrees; full smooth endpoint spaces require Fourier-series convergence/completion conventions rather than a naive algebraic direct sum of arbitrary modes.

For the standard sphere $S^{2n+1}=S(\mathcal O(-1))\to\PP^n$, the aid's complete mode count is
\[
 \dim H_b^{0,q}(S^{2n+1})_m=
 \begin{cases}
 \binom{m+n}{n},&q=0,\ m\ge0,\\
 \binom{-m-1}{n},&q=n,\ m\le-n-1,\\
 0,&\text{otherwise}.
 \end{cases}
\]
This fixes the weight sign and exhibits the infinite positive and negative endpoint tails. No conclusion about the unspecified Akahori problems is inferred from this special regular model.
\subsection{Short English statement}
The source records a topic discussed by Akahori but no recoverable mathematical question. The supplied auxiliary example identifies Fourier modes of Kohn--Rossi cohomology on a positive circle bundle with line-bundle Dolbeault cohomology. Interior degrees have only finitely many contributing modes, while sphere endpoint degrees have infinite tails.
\subsection{Sources}
\begin{itemize}
\item Local UnsolvedMath registry, permanent ID 20002090, source code catalogue entry; canonical \texttt{aim-geometry-notes.json}, notes[427], statement status reconstructed unverified.
\item AIM \emph{Holomorphic curves in contact geometry}, Question 1.32, discussion attributed to Akahori: \url{https://aimath.org/WWN/contactgeom2/contactgeom2.pdf}.
\item No separate authoritative citation for a precise Akahori problem or for the Fourier--Dolbeault special case is supplied. The aid invokes standard line-bundle cohomology, Serre vanishing and duality without bibliographic entries.
\item Registry research attempt 1, dated 2026-08-14: context only, status at run not a problem, rigor ``proved'', novelty confidence low. The machine-generated, non-peer-reviewed aid requires independent verification before citation.
\end{itemize}
\end{document}
